\documentclass{article}
\usepackage{amsmath}
\usepackage{amssymb}
\usepackage{enumitem}

\begin{document}

\title{Lecture Scribe \\ RANDOM VARIABLES, DISTRIBUTIONS \& EXPECTATION}
\maketitle

\section{(1) Previous Lecture Recap: Random Variables (RVs)}

\subsection{1.1 Definition of a Random Variable}
Let: \\
SSS = Sample Space of a random experiment \\
$\omega\in S$ = an outcome \\
A Random Variable (RV) is defined as a function: \\
$X:S\rightarrow \mathbb{R}$ \\
This means: \\
For every outcome $\omega\in S$, the function assigns a real number. \\
So, \\
$X(\omega)=$ real number \\
Thus, a random variable converts outcomes of an experiment into numerical values.

\subsection{1.2 Types of Random Variables}

\subsubsection{(A) Discrete Random Variable}
A random variable that takes countable values. \\
Representation: \\
$X=x_1,x_2,x_3,\dots$

\subsubsection{(B) Continuous Random Variable}
A random variable that takes values over a continuous interval. \\
Representation: \\
$a\le X\le b$

\subsection{1.3 Probability Mass Function (PMF)}
If $X$ is discrete: \\
$P(X=x_i)=p_i$ \\
Properties \\
$0\le p_i\le 1$ \\
$\sum p_i=1$ \\
These conditions ensure that probabilities are valid and total probability equals 1.

\section{(2) Independent Events}

\subsection{2.1 Definition}
Two events $A$ and $B$ are independent if: \\
$P(A \cap B) = P(A)P(B)$ \\
Logical meaning: \\
The occurrence of one event does not change the probability of the other event.

\subsection{2.2 Example: Two Coin Tosses}
Step 1: Write Sample Space \\
$S=\{HH, HT, TH, TT\}$ \\
Total outcomes = 4.

Step 2: Define Events \\
Let: \\
$A=$ First toss is Head \\
Outcomes satisfying $A$: \\
$\{HH, HT\}$ \\
Thus, \\
$P(A) = \frac{2}{4} = \frac{1}{2}$

Let: \\
$B=$ Second toss is Head \\
Outcomes satisfying $B$: \\
$\{HH, TH\}$ \\
Thus, \\
$P(B) = \frac{2}{4} = \frac{1}{2}$

Step 3: Find Intersection \\
$A \cap B = \{HH\}$ \\
Thus, \\
$P(A \cap B) = \frac{1}{4}$

Step 4: Verify Independence Condition \\
$P(A)P(B) = \frac{1}{2} \times \frac{1}{2} = \frac{1}{4}$ \\
Since: \\
$P(A \cap B) = P(A)P(B)$ \\
Therefore, $A$ and $B$ are independent.

\section{(3) Types of Discrete Random Variables}

\subsection{3.1 Bernoulli Random Variable}

\subsubsection{Definition}
A random variable with two possible outcomes: \\
$X = \begin{cases} 1 & \text{with probability } p \\ 0 & \text{with probability } 1-p \end{cases}$

\subsubsection{PMF}
$P(X=x) = p^x (1-p)^{1-x}$ \\
where $x=0,1$

\subsubsection{Expectation (Step-by-Step)}
By definition: \\
$E[X] = \sum xP(X=x)$ \\
Substitute possible values: \\
$E[X] = 1 \cdot p + 0 \cdot (1-p)$ \\
Compute terms: \\
$1 \cdot p = p$ \quad $0 \cdot (1-p) = 0$ \\
Add: \\
$E[X] = p + 0 = p$

\subsection{3.2 Binomial Random Variable}

\subsubsection{Definition}
If: \\
$n$ independent trials \\
Each trial has probability of success $p$ \\
Then: \\
$X=$ Number of successes

\subsubsection{PMF}
$P(X=k) = \binom{n}{k} p^k (1-p)^{n-k}$

\subsubsection{Expectation (Step-by-Step Logical Derivation)}
Let: \\
$X = X_1 + X_2 + \dots + X_n$ \\
Each $X_i$ is Bernoulli. \\
We know: \\
$E[X_i] = p$ \\
Now apply linearity: \\
$E[X] = E[X_1 + X_2 + \dots + X_n] = E[X_1] + E[X_2] + \dots + E[X_n]$ \\
Substitute: \\
$= p + p + \dots + p$ \\
Since there are $n$ terms: \\
$E[X] = np$

\subsection{3.3 Geometric Random Variable}

\subsubsection{Definition}
Number of trials required to get the first success.

\subsubsection{PMF}
$P(X=k) = (1-p)^{k-1}p$

\subsubsection{Expectation (Step-by-Step Derivation)}
$E[X] = \sum_{k=1}^{\infty} k(1-p)^{k-1}p$ \\
Factor $p$: \\
$= p \sum_{k=1}^{\infty} k(1-p)^{k-1}$ \\
Let: \\
$r = 1-p$ \\
We use the known series result: \\
$\sum_{k=1}^{\infty} kr^{k-1} = \frac{1}{(1-r)^2}$ \\
Substitute: \\
$1-r = p$ \\
Thus: \\
$\sum k(1-p)^{k-1} = \frac{1}{p^2}$ \\
Therefore: \\
$E[X] = p \cdot \frac{1}{p^2} = \frac{1}{p}$

\subsection{3.4 Poisson Random Variable}

\subsubsection{Definition}
Used when events occur independently in fixed interval. \\
Parameter: \\
$\lambda > 0$

\subsubsection{PMF}
$P(X=k) = \frac{e^{-\lambda}\lambda^k}{k!}$

\subsubsection{Expectation (Step-by-Step)}
$E[X] = \sum_{k=0}^{\infty} k\frac{e^{-\lambda}\lambda^k}{k!}$ \\
Rewrite: \\
$\frac{k}{k!} = \frac{1}{(k-1)!}$ \\
Thus: \\
$E[X] = \lambda e^{-\lambda} \sum_{k=1}^{\infty} \frac{\lambda^{k-1}}{(k-1)!}$ \\
Let $m = k-1$: \\
$= \lambda e^{-\lambda} \sum_{m=0}^{\infty} \frac{\lambda^m}{m!}$ \\
Recognize exponential expansion: \\
$\sum_{m=0}^{\infty} \frac{\lambda^m}{m!} = e^\lambda$ \\
Thus: \\
$E[X] = \lambda e^{-\lambda} e^\lambda = \lambda$

\section{(4) Expectation of Random Variables}

\subsection{4.1 Definition}
$E[X] = \sum xP(X=x)$

\subsection{4.2 Expectation of Function}
$E[g(X)] = \sum g(x)P(X=x)$

\subsubsection{Example}
Given: \\
$P(X=1)=0.5, P(X=2)=0.5$ \\
Find: \\
$E[X^2] = 1^2(0.5) + 2^2(0.5) = 1(0.5) + 4(0.5) = 0.5 + 2 = 2.5$

\subsection{4.3 Linearity of Expectation}
$E[aX+b] = \sum (ax+b)P(X=x) = a\sum xP(X=x) + b\sum P(X=x) = aE[X] + b(1) = aE[X]+b$

\subsection{4.4 Variance Derivation}
$\text{Var}(X) = E[(X-\mu)^2]$ \\
Expand: \\
$(X-\mu)^2 = X^2 - 2\mu X + \mu^2$ \\
Take expectation: \\
$E[X^2] - 2\mu E[X] + \mu^2$ \\
Since $E[X]=\mu$: \\
$= E[X^2] - 2\mu^2 + \mu^2 = E[X^2] - \mu^2$ \\
Thus: \\
$\text{Var}(X) = E[X^2]-(E[X])^2$

\subsection{4.5 nᵗʰ Moments}
$\mu_n = E[X^n]$ \\
Special cases: \\
2nd central moment $\rightarrow$ Variance \\
3rd $\rightarrow$ Skewness \\
4th $\rightarrow$ Kurtosis

\section{(5) Cumulative Distribution Function (CDF)}

\subsection{Definition}
$F(x)=P(X\le x)$

\subsection{Properties}
$0\le F(x)\le 1$ \\
Non-decreasing \\
Right continuous \\
$\lim_{x\to-\infty}F(x)=0$ \\
$\lim_{x\to\infty}F(x)=1$

\subsection{Example}
Given: \\
$P(X=1)=0.4, P(X=2)=0.6$ \\
Then: \\
For $x<1$: \\
$F(x)=0$ \\
For $1\le x<2$: \\
$F(x)=0.4$ \\
For $x\ge 2$: \\
$F(x)=1$

\section{(6) Probability Density Function (PDF)}

\subsection{Definition}
$f(x)\ge 0$ \quad $\int_{-\infty}^{\infty} f(x)dx=1$

\subsection{Probability Between Two Points}
$P(a\le X\le b)=\int_a^b f(x)dx$

\subsection{Relation Between CDF and PDF}
$F(x)=\int_{-\infty}^x f(t)dt$ \\
Differentiate: \\
$f(x)=\frac{d}{dx}F(x)$

\section{FINAL REVISION SUMMARY}
\begin{itemize}
\item RV = Function from sample space to real numbers
\item Independence = $P(A\cap B)=P(A)P(B)$
\item Bernoulli $\rightarrow$ $E[X]=p$
\item Binomial $\rightarrow$ $E[X]=np$
\item Geometric $\rightarrow$ $E[X]=1/p$
\item Poisson $\rightarrow$ $E[X]=\lambda$
\item Variance $\rightarrow$ $E[X^2]-(E[X])^2$
\item CDF $\rightarrow$ $F(x)=P(X\le x)$
\item PDF integrates to 1
\end{itemize}

\end{document}
